% TOP_PROBLEM: {"problemId": "problem.alperin-mckay-conjecture", "canonicalTitle": "Alperin–McKay Conjecture", "releaseRank": 278, "primaryDomain": "algebra_representation_category", "primaryDomainLabel": "Algebra, representation & category theory", "displayStatus": "open_with_solved_subcases", "releaseStatus": "open_with_solved_subcases"}
% !TeX program = lualatex
% Definition and sources only; this file is not an LLM solution attempt.
\documentclass[11pt]{article}
\usepackage[margin=25mm]{geometry}
\usepackage{fontspec}
\setmainfont{Latin Modern Roman}
\usepackage{amsmath,amssymb,mathtools}
\usepackage{unicode-math}
\setmathfont{Latin Modern Math}
\usepackage{xurl,hyperref}
\hypersetup{colorlinks=true,urlcolor=blue}
\setcounter{secnumdepth}{0}
\setlength{\parindent}{0pt}
\setlength{\parskip}{5pt}
\setlength{\emergencystretch}{3em}
\begin{document}
\section*{\texorpdfstring{Alperin–McKay Conjecture}{Alperin–McKay Conjecture}}\label{278}
\subsection{Definitions and mathematical statement}
A $p$-block groups ordinary irreducible characters according to a primitive central idempotent in modular representation theory; its defect group $D$ is the associated maximal local $p$-subgroup, defined up to conjugacy. Let $b$ be the Brauer correspondent of a block $B$ in $N_G(D)$. For a positive integer $a$, write $a_p$ for its largest dividing power of $p$. Define
\[
 k_0(B)=\#\{\chi\in\operatorname{Irr}(B):\chi(1)_p=|G:D|_p\}.
\]
Define $k_0(b)$ by the same formula with ambient group $N_G(D)$. The Alperin--McKay assertion is $k_0(B)=k_0(b)$ for every finite $G$, every prime, and every block. It compares height-zero characters only, not all characters of the two blocks.
\par\smallskip\textit{Formulation and concept references:} \href{https://doi.org/10.1017/fms.2022.36}{[S1]}.

\subsection{Short English statement}
Does each finite-group block contain exactly as many height-zero ordinary characters as its Brauer correspondent in the normalizer of a defect group?
\subsection{Sources}
\begin{itemize}
\item[S1]Lucas Ruhstorfer, Quasi-isolated blocks and the Alperin–McKay conjecture, introduction; height-zero convention from The Alperin–McKay conjecture for the prime2, introduction.
\url{https://doi.org/10.1017/fms.2022.36}.
\textit{Formulation source.}
\end{itemize}
\end{document}
